\section{Turing Machine --- Build --- NDClock Host}
%%% generated by the blueprint pipeline from TCSlib.Complexity.TuringMachine.Build.NDClockHost
\subsection*{Turing Machine --- Build --- NDClock Host: Introduction}

Declarations of the Lean module \texttt{TCSlib.Complexity.TuringMachine.Build.NDClockHost}.

\subsection*{Turing Machine --- Build --- NDClock Host: Results}

\begin{definition}[Capped output update]
\lean{Turing.capPush}
\label{Turing.capPush}
\leanok
For a current cap value $c\in\{\text{none},\text{a Boolean symbol},\text{overflow}\}$
and an optional emitted Boolean $e$, the updated cap is unchanged if nothing is emitted;
if a symbol is emitted and the cap was empty, it records that symbol; if a symbol is
emitted when the cap was already nonempty, it records overflow.
\end{definition}

\begin{definition}[Cap of a Boolean list]
\lean{Turing.capOf}
\label{Turing.capOf}
\leanok
For a finite list of Boolean symbols, its cap is empty when the list is empty, records
its symbol when the list has exactly one symbol, and records overflow when the list has
at least two symbols.
\end{definition}

\begin{theorem}[Cap of an Appended Bit]
\lean{Turing.capOf_append}
\label{Turing.capOf_append}
\leanok
\uses{Turing.capOf, Turing.capPush}
\difficulty{4}
For every Boolean word $o$ and optional Boolean $e$, the cap of the word formed by
appending $e$ (if present) to $o$ equals the cap-push operation applied to the cap of
$o$ and $e$.
\end{theorem}

\begin{theorem}[Singleton Characterization of the Cap]
\lean{Turing.capOf_eq_some_some}
\label{Turing.capOf_eq_some_some}
\leanok
\uses{Turing.capOf}
\difficulty{4}
For every Boolean word $o$ and Boolean $v$, the cap of $o$ equals the nested optional
value containing $v$ exactly when $o$ is the singleton word $[v]$.
\end{theorem}

\begin{definition}[Phases of a clocked simulation]
\lean{Turing.ClockPhase}
\label{Turing.ClockPhase}
\leanok
For state types $SS$ and $AS$, the clock phase is a boot phase carrying a yardstick
machine state; one of two rewind phases; a run phase carrying an optional attempt state
and a captured value of type $\mathrm{Option}(\mathrm{Option}(\mathrm{Bool}))$; or one
of three reset-sweep phases, in leftward, rightward, then leftward order.
\end{definition}

\begin{definition}[Combined tape slots]
\lean{Turing.ClockHost.slots}
\label{Turing.ClockHost.slots}
\leanok
Given a type $\alpha$, nonnegative integers $s,a$, functions
$f:\mathrm{Fin}(s)\to\alpha$ and $g:\mathrm{Fin}(a)\to\alpha$, and an element
$y\in\alpha$, define a function on $s+a+1$ slots by placing the values of $f$ first,
then the values of $g$, and then the constant value $y$ in the final slot.
\end{definition}

\begin{definition}[State-machine slot indices]
\lean{Turing.ClockHost.sI}
\label{Turing.ClockHost.sI}
\leanok
For nonnegative integers $s,a$ and $i\in\mathrm{Fin}(s)$, the associated index is the
corresponding position among the first $s$ slots of a list of $s+a+1$ slots.
\end{definition}

\begin{definition}[Attempt-machine slot indices]
\lean{Turing.ClockHost.aI}
\label{Turing.ClockHost.aI}
\leanok
For nonnegative integers $s,a$ and $i\in\mathrm{Fin}(a)$, the associated index is
position $s+i$ among the $a$ slots following the first $s$ slots of a list of $s+a+1$
slots.
\end{definition}

\begin{definition}[Final slot index]
\lean{Turing.ClockHost.yI}
\label{Turing.ClockHost.yI}
\leanok
For nonnegative integers $s$ and $a$, let the final slot index be the last position in a
list of $s+a+1$ slots.
\end{definition}

\begin{theorem}[Assignment to the first tape block]
\lean{Turing.ClockHost.slots_sI}
\label{Turing.ClockHost.slots_sI}
\leanok
\uses{Turing.ClockHost.sI, Turing.ClockHost.slots}
\difficulty{1}
Let $s,a$ be natural numbers, let $f$ and $g$ be functions on the first and second
blocks of a combined index set of size $s+a+1$, and let $y$ be an element of their
codomain. At the combined index corresponding to any index $i$ in the first block, the
function assembled from $f$, $g$, and $y$ has value $f(i)$.
\end{theorem}

\begin{theorem}[Assignment to the second tape block]
\lean{Turing.ClockHost.slots_aI}
\label{Turing.ClockHost.slots_aI}
\leanok
\uses{Turing.ClockHost.aI, Turing.ClockHost.slots}
\difficulty{1}
Let $s,a$ be natural numbers, let $f$ and $g$ be functions on the first and second
blocks of a combined index set of size $s+a+1$, and let $y$ be an element of their
codomain. At the combined index corresponding to any index $i$ in the second block, the
function assembled from $f$, $g$, and $y$ has value $g(i)$.
\end{theorem}

\begin{theorem}[Assignment to the final tape slot]
\lean{Turing.ClockHost.slots_yI}
\label{Turing.ClockHost.slots_yI}
\leanok
\uses{Turing.ClockHost.slots, Turing.ClockHost.yI}
\difficulty{1}
Let $s,a$ be natural numbers, let $f$ and $g$ be functions on the first and second
blocks of a combined index set of size $s+a+1$, and let $y$ be an element of their
codomain. At the final, singleton index, the function assembled from $f$, $g$, and $y$
has value $y$.
\end{theorem}

\begin{theorem}[Exhaustion of the combined indices]
\lean{Turing.ClockHost.slots_cases}
\label{Turing.ClockHost.slots_cases}
\leanok
\uses{Turing.ClockHost.aI, Turing.ClockHost.sI, Turing.ClockHost.yI}
\difficulty{4}
Let $s,a$ be natural numbers and let $P$ be a property of the indices in a set of size
$s+a+1$. Suppose $P$ holds at every index in the first block, at every index in the
second block, and at the single remaining index. Then $P$ holds at every index in the
combined set.
\end{theorem}

\begin{definition}[Work tape after an action]
\lean{Turing.ClockHost.tapeAfter}
\label{Turing.ClockHost.tapeAfter}
\leanok
Given an action consisting of a captured output value and a head movement, a tape
$t:\mathbb{Z}\to\mathrm{Option}(\mathrm{Bool})$, and a head position $p\in\mathbb{Z}$,
define the resulting tape by applying the action to a one-tape configuration with absent
state, empty output, and sole work tape and head position $t$ and $p$.
\end{definition}

\begin{definition}[Clocked host configuration]
\lean{Turing.ClockHost.hostCfg}
\label{Turing.ClockHost.hostCfg}
\leanok
\uses{Turing.Cfg, Turing.ClockHost.slots, Turing.ClockPhase}
Let $SS$ and $AS$ be state types, let $s,a$ be nonnegative integers, and let $x$ be a
finite Boolean input. Given an optional clock phase, an input-head position for $x$,
state-machine tapes and head positions, attempt-machine tapes and head positions, a
final tape and head position, and an output string, define the host configuration to
have clock-phase state, those input and output data, and work tapes and head positions
formed by concatenating the state-machine slots, the attempt-machine slots, and the
final slot.
\end{definition}

\begin{theorem}[Applying an Action to a Clock-Host Configuration]
\lean{Turing.ClockHost.apply_hostCfg}
\label{Turing.ClockHost.apply_hostCfg}
\leanok
\uses{Turing.Action, Turing.Action.apply, Turing.ClockHost.hostCfg, Turing.ClockHost.slots, Turing.ClockHost.slots_aI, Turing.ClockHost.slots_cases, Turing.ClockHost.slots_sI, Turing.ClockHost.slots_yI, Turing.ClockHost.tapeAfter, Turing.ClockPhase, Turing.moveInputPos}
\difficulty{5}
Let $s,a$ be nonnegative integers, let $x$ be a Boolean input word, and choose an
input-head movement, one write-and-move instruction for each of the $s$ yardstick tapes,
one for each of the $a$ simulated tapes, and one for the clock tape. Also choose an
optional output bit, initial and successor clock-phase states, an input-head position,
the contents and head positions of each group of tapes, and an output word. Applying the
assembled action to the host configuration with initial phase and these tape data yields
the host configuration with successor phase, the input head moved by the chosen clamped
movement, each tape updated by its instruction with its head advanced by the chosen
movement, and the output bit appended to the output word.
\end{theorem}

\begin{theorem}[Work Symbols on Yardstick Tapes]
\lean{Turing.ClockHost.hostCfg_symbols_sI}
\label{Turing.ClockHost.hostCfg_symbols_sI}
\leanok
\uses{Turing.Cfg.workTapeSymbols, Turing.ClockHost.hostCfg, Turing.ClockHost.sI, Turing.ClockPhase}
\difficulty{1}
Let $s$ and $a$ be natural numbers, let $x$ be a Boolean input word, and choose a clock
phase, an input-head position, tape contents and head positions for the $s$ yardstick
tapes, the $a$ attempt tapes, and the extra tape, together with an output word. For
every yardstick-tape index $i$, the symbol read by the head at that index in the
resulting host configuration is the content of yardstick tape $i$ at its chosen head
position.
\end{theorem}

\begin{theorem}[Work Symbols on Attempt Tapes]
\lean{Turing.ClockHost.hostCfg_symbols_aI}
\label{Turing.ClockHost.hostCfg_symbols_aI}
\leanok
\uses{Turing.Cfg.workTapeSymbols, Turing.ClockHost.aI, Turing.ClockHost.hostCfg, Turing.ClockPhase}
\difficulty{1}
Let $s$ and $a$ be natural numbers, let $x$ be a Boolean input word, and choose a clock
phase, an input-head position, tape contents and head positions for the $s$ yardstick
tapes, the $a$ attempt tapes, and the extra tape, together with an output word. For
every attempt-tape index $i$, the symbol read by the head at that index in the resulting
host configuration is the content of attempt tape $i$ at its chosen head position.
\end{theorem}

\begin{theorem}[Symbol on the Extra Tape]
\lean{Turing.ClockHost.hostCfg_symbols_yI}
\label{Turing.ClockHost.hostCfg_symbols_yI}
\leanok
\uses{Turing.Cfg.workTapeSymbols, Turing.ClockHost.hostCfg, Turing.ClockHost.yI, Turing.ClockPhase}
\difficulty{1}
Let $s$ and $a$ be natural numbers, let $x$ be a Boolean input word, and choose a clock
phase, an input-head position, tape contents and head positions for the $s$ yardstick
tapes, the $a$ attempt tapes, and the extra tape, together with an output word. The
symbol read by the head on the extra tape in the resulting host configuration is the
content of that tape at its chosen head position.
\end{theorem}

\begin{definition}[Idle action]
\lean{Turing.ClockHost.idle}
\label{Turing.ClockHost.idle}
\leanok
For a Boolean finite-state Turing machine $S$, a Boolean nondeterministic Turing machine
$A$, and a Boolean parameter $loop$, define the idle action to consist of no captured
tape value and zero head movement.
\end{definition}

\begin{theorem}[Idle Tape Heads]
\lean{Turing.ClockHost.heads_idle}
\label{Turing.ClockHost.heads_idle}
\leanok
\uses{Turing.ClockHost.idle}
\difficulty{3}
For every natural number $n$ and every function $P:	ext{Fin}(n)\to\mathbb{Z}$ assigning
integer head positions, adding the movement specified by the idle action leaves $P$
unchanged.
\end{theorem}

\begin{definition}[Output capture action]
\lean{Turing.ClockHost.emitTo}
\label{Turing.ClockHost.emitTo}
\leanok
For a deterministic Boolean machine $S$, a nondeterministic Boolean machine $A$, and a
Boolean loop flag, define the output action associated to an optional Boolean output $e$
as follows: if $e$ is present, write it to the output and move the associated head
right; if it is absent, make no write and leave that head stationary.
\end{definition}

\begin{definition}[Clocked transition function]
\lean{Turing.ClockHost.clockTr}
\label{Turing.ClockHost.clockTr}
\leanok
\uses{Turing.Action, Turing.ClockHost.aI, Turing.ClockHost.emitTo, Turing.ClockHost.idle, Turing.ClockHost.sI, Turing.ClockHost.slots, Turing.ClockHost.yI, Turing.ClockPhase, Turing.capPush}
Given a deterministic Boolean machine $S$, a nondeterministic Boolean machine $A$, a
Boolean loop flag, a choice bit, a clock phase, the possibly blank input symbol under
the head, and the symbols under the host machine's work-tape heads, the transition
returns the following action. In the boot phase it simulates one step of $S$, routes its
work-tape actions to the corresponding tapes, and enters the next boot phase or begins
rewinding when $S$ halts. It then rewinds the yardstick tape and starts $A$; during the
run phase it simulates the chosen transition of $A$ while the yardstick remains,
captures its output, and when the yardstick is exhausted either emits the captured
Boolean output if $A$ halted with one, halts if the loop flag is false, or performs the
three reset sweeps and restarts $A$ if the loop flag is true. The rewind and reset
phases move the relevant heads as specified by the tape contents, and all other tape
actions and outputs are blank or absent.
\end{definition}

\begin{definition}[Clock host machine]
\lean{Turing.ClockHost.hostTM}
\label{Turing.ClockHost.hostTM}
\leanok
\uses{Turing.ClockHost.clockTr, Turing.ClockPhase, Turing.NDTM}
Given a deterministic Boolean Turing machine $S$, a nondeterministic Boolean Turing
machine $A$, and a Boolean parameter $	extit{loop}$, define a nondeterministic Turing
machine with $S.k + A.k + 1$ work tapes and states given by the clock phases associated
with the state sets of $S$ and $A$. Its initial state is the boot phase containing the
initial state of $S$, and its transition function is the clocked transition rule
determined by $S$, $A$, and $	extit{loop}$.
\end{definition}

\begin{definition}[Finite clock-host machine]
\lean{Turing.FinNDTM.clockHost}
\label{Turing.FinNDTM.clockHost}
\leanok
\uses{Turing.ClockHost.hostTM, Turing.ClockPhase, Turing.FinNDTM, Turing.FinTM}
Given a finite deterministic Turing machine $S$ over the Boolean alphabet, a finite
nondeterministic Turing machine $A$ over the Boolean alphabet, and a Boolean parameter
specifying whether to loop, define a finite nondeterministic Turing machine with
$S.k+A.k+1$ work tapes. Its states are clock phases associated with the states of $S$
and $A$, and its transitions are those of the clock-host construction from $S$, $A$, and
the loop parameter.
\end{definition}

\begin{theorem}[Host step from a host configuration]
\lean{Turing.ClockHost.stepWith_hostCfg}
\label{Turing.ClockHost.stepWith_hostCfg}
\leanok
\uses{Turing.Action.apply, Turing.Cfg.inputSymbol, Turing.Cfg.workTapeSymbols, Turing.ClockHost.clockTr, Turing.ClockHost.hostCfg, Turing.ClockHost.hostTM, Turing.ClockPhase, Turing.NDTM.stepWith, Turing.NDTM.stepWith_of_state}
\difficulty{1}
Let $S$ be a finite-state deterministic Boolean machine, let $A$ be a nondeterministic
Boolean machine, and fix a Boolean loop parameter and Boolean input $x$. For any Boolean
choice, clock phase, input-head position, work-tape contents and head positions for both
machines and the clock tape, and output word, taking one host-machine step from the host
configuration with that data gives the same configuration as applying the action
prescribed by the host transition function at that phase, choice, scanned input symbol,
and scanned work-tape symbols.
\end{theorem}

\begin{definition}[Yardstick tape contents]
\lean{Turing.ClockHost.yard}
\label{Turing.ClockHost.yard}
\leanok
\uses{Turing.FinTM.bufferTape}
For each natural number $m$, define the yardstick tape to contain $m$ consecutive true
symbols starting at the origin, with blank cells elsewhere.
\end{definition}

\begin{theorem}[Contents of the yardstick tape]
\lean{Turing.ClockHost.yard_of_lt}
\label{Turing.ClockHost.yard_of_lt}
\leanok
\uses{Turing.ClockHost.yard, Turing.FinTM.bufferTape_nat}
\difficulty{3}
For every natural number $m$ and integer cell $z$, if $0 \le z$ and $z < m$, then the
yardstick tape of length $m$ contains the bit $1$ at cell $z$.
\end{theorem}

\begin{theorem}[Negative yardstick cell is blank]
\lean{Turing.ClockHost.yard_neg_one}
\label{Turing.ClockHost.yard_neg_one}
\leanok
\uses{Turing.ClockHost.yard, Turing.FinTM.bufferTape_left}
\difficulty{1}
For every natural number $m$, the yardstick tape of length $m$ is blank at cell $-1$.
\end{theorem}

\begin{theorem}[Yardstick endpoint is blank]
\lean{Turing.ClockHost.yard_len}
\label{Turing.ClockHost.yard_len}
\leanok
\uses{Turing.ClockHost.yard, Turing.FinTM.bufferTape_nat}
\difficulty{1}
For every natural number $m$, the yardstick tape of length $m$ is blank at cell $m$.
\end{theorem}

\begin{definition}[Boot configuration]
\lean{Turing.ClockHost.bootCfg}
\label{Turing.ClockHost.bootCfg}
\leanok
\uses{Turing.Cfg, Turing.ClockHost.hostCfg, Turing.ClockPhase, Turing.FinTM.bufferTape}
Let $S$ be a finite deterministic Boolean machine, $A$ a finite nondeterministic Boolean
machine, and let the loop flag be Boolean. Given a configuration of $S$ on input $x$,
define the host configuration with the same input-head position, $S$'s work-tape
contents and head positions, and output; its state is the boot phase with $S$'s current
state when present, or the rewind phase when $S$ has halted. The additional work tapes
are blank, with their heads at the origin.
\end{definition}

\begin{theorem}[Boot Step Simulation]
\lean{Turing.ClockHost.boot_step}
\label{Turing.ClockHost.boot_step}
\leanok
\uses{Turing.Action.apply, Turing.Cfg, Turing.Cfg.inputSymbol, Turing.Cfg.workTapeSymbols, Turing.ClockHost.apply_hostCfg, Turing.ClockHost.bootCfg, Turing.ClockHost.clockTr, Turing.ClockHost.emitTo, Turing.ClockHost.hostCfg, Turing.ClockHost.hostTM, Turing.ClockHost.sI, Turing.ClockHost.tapeAfter, Turing.FinTM.bufferTape_append, Turing.MultiTapeTM.step, Turing.NDTM.stepWith, Turing.NDTM.stepWith_of_state}
\difficulty{5}
Let $S$ be a deterministic Boolean Turing machine, $A$ a nondeterministic Boolean Turing
machine, and let the clock loop flag and choice bit $b$ be fixed. For a configuration
$c$ of $S$ on a Boolean input, suppose its state is $q$. Then one host-machine step
under $b$, starting from the boot configuration of $c$, equals the boot configuration of
the one-step successor of $c$ under $S$.
\end{theorem}

\begin{theorem}[Boot Run Simulation]
\lean{Turing.ClockHost.boot_run}
\label{Turing.ClockHost.boot_run}
\leanok
\uses{Turing.Cfg, Turing.ClockHost.bootCfg, Turing.ClockHost.boot_step, Turing.ClockHost.hostTM, Turing.MultiTapeTM.runFrom, Turing.MultiTapeTM.runFrom_succ_eq_step, Turing.MultiTapeTM.runFrom_zero, Turing.NDTM.runWith, Turing.NDTM.runWith_cons}
\difficulty{4}
Let $S$ be a deterministic Boolean Turing machine, $A$ a nondeterministic Boolean Turing
machine, and let the clock loop flag be fixed. For a configuration $c$ of $S$ on a
Boolean input and a finite Boolean choice word $w$, suppose $S$ is still running after
every $t<|w|$ steps from $c$. Then the host machine run from the boot configuration of
$c$ under $w$ equals the boot configuration of the configuration reached by running $S$
from $c$ for $|w|$ steps.
\end{theorem}

\begin{theorem}[Initial host and boot configurations]
\lean{Turing.ClockHost.initCfg_eq_bootCfg}
\label{Turing.ClockHost.initCfg_eq_bootCfg}
\leanok
\uses{Turing.Cfg.init, Turing.ClockHost.bootCfg, Turing.ClockHost.hostCfg, Turing.ClockHost.hostTM, Turing.ClockHost.slots_cases, Turing.MultiTapeTM.initCfg, Turing.NDTM.initCfg}
\difficulty{4}
For a deterministic machine $S$, a nondeterministic machine $A$, a Boolean loop flag,
and a Boolean input word $x$, the initial configuration of the host machine on $x$
equals the boot configuration formed from the initial configuration of $S$ on $x$.
\end{theorem}

\begin{definition}[Repeated leftward input movement]
\lean{Turing.ClockHost.inLeft}
\label{Turing.ClockHost.inLeft}
\leanok
\uses{Turing.moveInputPos}
For a natural number $n$, an input-head position among the $n+2$ valid positions, and a
natural number $j$, define the position obtained by attempting to move left $j$ times,
using clamped movement at the input boundaries.
\end{definition}

\begin{theorem}[Position after repeated left moves]
\lean{Turing.ClockHost.inLeft_val}
\label{Turing.ClockHost.inLeft_val}
\leanok
\uses{Turing.ClockHost.inLeft, Turing.moveInputPos_neg_of_ne_left}
\difficulty{4}
For every $n\in\mathbb N$, every input-head position $ip$ among the $n+2$ positions, and
every $j\in\mathbb N$, the position reached after applying the clamped left move $j$
times has value $\max(ip- j,0)$, where subtraction is truncated in $\mathbb N$.
\end{theorem}

\begin{theorem}[Right move after reaching the left boundary]
\lean{Turing.ClockHost.moveInputPos_inLeft}
\label{Turing.ClockHost.moveInputPos_inLeft}
\leanok
\uses{Turing.ClockHost.inLeft, Turing.ClockHost.inLeft_val, Turing.moveInputPos, Turing.moveInputPos_pos_of_ne_right}
\difficulty{3}
For every $n\in\mathbb N$, input-head position $ip$ among the $n+2$ positions, and
$j\in\mathbb N$ with $ip\le j$, if the head is moved left $j$ times with clamping, then
moving it right once gives position $1$.
\end{theorem}

\begin{theorem}[A Rewind Step on a Marked Cell]
\lean{Turing.ClockHost.rew2_step}
\label{Turing.ClockHost.rew2_step}
\leanok
\uses{Turing.ClockHost.apply_hostCfg, Turing.ClockHost.clockTr, Turing.ClockHost.hostCfg, Turing.ClockHost.hostCfg_symbols_yI, Turing.ClockHost.hostTM, Turing.ClockHost.stepWith_hostCfg, Turing.ClockHost.yard, Turing.NDTM.stepWith, Turing.moveInputPos}
\difficulty{3}
Let $S$ be a deterministic Boolean Turing machine, $A$ a nondeterministic Boolean Turing
machine, $x$ a Boolean input, and $m$ a natural number. If the yard tape containing $m$
consecutive true symbols has value true at an integer position $z$, then one step of the
host machine in phase $	ext{rew2}$ moves the input head one position left and the yard
head from $z$ to $z-1$, leaving the phase, all other tape contents and head positions,
and the output unchanged.
\end{theorem}

\begin{theorem}[First rewind step]
\lean{Turing.ClockHost.rew_step}
\label{Turing.ClockHost.rew_step}
\leanok
\uses{Turing.ClockHost.apply_hostCfg, Turing.ClockHost.clockTr, Turing.ClockHost.hostCfg, Turing.ClockHost.hostTM, Turing.ClockHost.stepWith_hostCfg, Turing.ClockHost.yard, Turing.NDTM.stepWith, Turing.moveInputPos}
\difficulty{3}
Let $S$ be a deterministic Boolean Turing machine with $S.k$ work tapes and $A$ a
nondeterministic Boolean Turing machine with $A.k$ work tapes. For every Boolean choice
$b$, Boolean input $x$, natural number $m$, input-head position $ip \in
\mathrm{Fin}(|x|+2)$, tape contents $sT,aT$ and head positions $sP,aP$ for the work
tapes of $S$ and $A$, and integer $z$, one step of the host machine from phase
$\mathrm{rew}$, with yard tape storing $m$ true symbols from cell zero and its head at
$z$, enters phase $\mathrm{rew2}$, moves the input head according to the clamped left
move, moves the yard head to $z-1$, and leaves the other work tapes and the output
unchanged.
\end{theorem}

\begin{theorem}[Transition from Rewind to Run]
\lean{Turing.ClockHost.rew2_turn}
\label{Turing.ClockHost.rew2_turn}
\leanok
\uses{Turing.ClockHost.apply_hostCfg, Turing.ClockHost.clockTr, Turing.ClockHost.hostCfg, Turing.ClockHost.hostCfg_symbols_yI, Turing.ClockHost.hostTM, Turing.ClockHost.stepWith_hostCfg, Turing.ClockHost.yard, Turing.ClockHost.yard_neg_one, Turing.NDTM.stepWith, Turing.moveInputPos}
\difficulty{3}
Let $S$ be a deterministic Boolean Turing machine and $A$ a nondeterministic Boolean
Turing machine, and fix a Boolean loop parameter, a Boolean input word $x$, a natural
number $m$, and a choice bit $b$. For any input-head position among the $|x|+2$
positions and any contents and head positions for the work tapes of $S$ and $A$, if the
yard tape stores a word of $m$ true symbols starting at cell $0$ and its head is at cell
$-1$, one step of the host machine enters the run phase with $A$ in its initial state
and no captured output. The input head moves one position right, the yard head moves to
cell $0$, and the other tape contents, their head positions, and the output word remain
unchanged.
\end{theorem}

\begin{theorem}[Rewinding the Yard Tape]
\lean{Turing.ClockHost.rew2_marks}
\label{Turing.ClockHost.rew2_marks}
\leanok
\uses{Turing.ClockHost.hostCfg, Turing.ClockHost.hostTM, Turing.ClockHost.inLeft, Turing.ClockHost.rew2_step, Turing.ClockHost.yard, Turing.ClockHost.yard_of_lt, Turing.NDTM.runWith, Turing.NDTM.runWith_cons}
\difficulty{4}
Let $S$ be a deterministic Boolean Turing machine, let $A$ be a nondeterministic Boolean
Turing machine, let $x$ be a finite Boolean input, let $m$ be a natural number, and fix
a Boolean loop setting. Choose an input-head position for $x$, work-tape contents and
head positions for the work tapes of $S$ and $A$, and a Boolean word $w$ with $|w| \le
m$. Starting the host machine in its second rewind phase, with the yard tape containing
$m$ consecutive marks and its head at position $|w|-1$, and with empty output, running
it on the choice word $w$ leaves it in the same phase with the yard-tape head at $-1$;
the yard tape contents and all work-tape contents and head positions are unchanged,
while the input head has moved left by $|w|$ positions, clamped at the input boundaries,
and the output remains empty.
\end{theorem}

\begin{theorem}[Running Through the Rewind Phase]
\lean{Turing.ClockHost.rewind_run}
\label{Turing.ClockHost.rewind_run}
\leanok
\uses{Turing.ClockHost.hostCfg, Turing.ClockHost.hostTM, Turing.ClockHost.inLeft, Turing.ClockHost.rew2_marks, Turing.ClockHost.rew2_turn, Turing.ClockHost.rew_step, Turing.ClockHost.yard, Turing.NDTM.runWith, Turing.NDTM.runWith_append, Turing.NDTM.runWith_cons, Turing.NDTM.runWith_nil, Turing.moveInputPos}
\difficulty{5}
Let $S$ be a deterministic Boolean Turing machine, $A$ a nondeterministic Boolean Turing
machine, $x$ a Boolean input, and $m$ a natural number. For any input-head position and
work-tape contents and head positions, if a choice word $w$ has length $m+2$, then
running the host machine from phase $	ext{rew}$ with the yard head at $m$ for $w$ steps
ends in phase $	ext{run}$ with $A$ in its initial state and no captured output, with the
yard head at $0$ and the input head at the position obtained by moving left $m+1$ steps
and then right once; the other tape data and output are unchanged.
\end{theorem}

\begin{definition}[Run configuration]
\lean{Turing.ClockHost.runCfg}
\label{Turing.ClockHost.runCfg}
\leanok
\uses{Turing.Cfg, Turing.ClockHost.hostCfg, Turing.ClockHost.yard, Turing.ClockPhase, Turing.capOf}
Let $S$ be a finite deterministic Boolean machine, $A$ a finite nondeterministic Boolean
machine, a Boolean loop flag, an input list $x$, and a natural number $m$ be given. From
a configuration of $A$ on $x$, an integer $j$, and contents and head positions for the
work tapes of $S$, define the host configuration whose state records $A$'s optional
current state and captured output, whose input-head position is $A$'s input-head
position, whose work tapes combine the supplied contents for $S$ with $A$'s
configuration, and whose yardstick tape contains $m$ consecutive true symbols starting
at the origin; the yardstick head is at $j$, and the output is empty.
\end{definition}

\begin{theorem}[Clocked simulation step]
\lean{Turing.ClockHost.run_step}
\label{Turing.ClockHost.run_step}
\leanok
\uses{Turing.Action.apply, Turing.Cfg, Turing.Cfg.inputSymbol, Turing.Cfg.workTapeSymbols, Turing.ClockHost.aI, Turing.ClockHost.apply_hostCfg, Turing.ClockHost.clockTr, Turing.ClockHost.hostCfg, Turing.ClockHost.hostCfg_symbols_yI, Turing.ClockHost.hostTM, Turing.ClockHost.runCfg, Turing.ClockHost.yI, Turing.ClockHost.yard_of_lt, Turing.NDTM.stepWith, Turing.NDTM.stepWith_of_halt, Turing.NDTM.stepWith_of_state, Turing.capOf, Turing.capOf_append}
\difficulty{6}
Let $S$ be a finite-state deterministic machine over the Boolean alphabet and let $A$ be
a nondeterministic machine over that alphabet, with a Boolean parameter indicating
whether to loop. For an input $x$, a configuration $a$ of $A$, a natural number $j<m$, a
Boolean choice $b$, and arbitrary work-tape contents and head positions for $S$, the
host machine’s step from the run configuration encoding $a$ at clock position $j$ equals
the run configuration encoding the result of taking one step of $A$ from $a$ using
choice $b$, with clock position $j+1$.
\end{theorem}

\begin{theorem}[Simulating a Run Within the Clock]
\lean{Turing.ClockHost.run_run}
\label{Turing.ClockHost.run_run}
\leanok
\uses{Turing.Cfg, Turing.ClockHost.hostTM, Turing.ClockHost.runCfg, Turing.ClockHost.run_step, Turing.NDTM.runWith, Turing.NDTM.runWith_cons}
\difficulty{4}
Let $S$ be a deterministic Boolean Turing machine, $A$ a nondeterministic Boolean Turing
machine, $x$ a Boolean input, and $m$ a natural number. Let $a$ be any configuration of
$A$ on $x$, let $j$ be a natural number, and let $u$ be a Boolean choice word such that
$j+|u|\le m$. For any work-tape contents and head positions of $S$, running the host
machine on $u$ from the representation of $a$ in phase $	ext{run}$, with yard-head
position $j$, gives the representation of the run of $A$ on $u$ from $a$, with yard-head
position $j+|u|$; the other stated host data are preserved as specified by that
representation.
\end{theorem}

\begin{theorem}[Run Simulation from Yard Position Zero]
\lean{Turing.ClockHost.run_run_zero}
\label{Turing.ClockHost.run_run_zero}
\leanok
\uses{Turing.Cfg, Turing.ClockHost.hostTM, Turing.ClockHost.runCfg, Turing.ClockHost.run_run, Turing.NDTM.runWith}
\difficulty{1}
Let $S$ be a deterministic Boolean Turing machine, $A$ a nondeterministic Boolean Turing
machine, $x$ a Boolean input, and $m$ a natural number. For any configuration $a$ of $A$
on $x$, any work-tape contents and head positions for $S$, and any Boolean choice word
$u$ with $|u|\le m$, running the host machine on $u$ from the representation of $a$ in
phase $	ext{run}$ with yard-head position zero gives the representation of the run of
$A$ on $u$ from $a$, with yard-head position $|u|$.
\end{theorem}

\begin{theorem}[Successful timeout step]
\lean{Turing.ClockHost.timeout_success}
\label{Turing.ClockHost.timeout_success}
\leanok
\uses{Turing.Cfg, Turing.ClockHost.apply_hostCfg, Turing.ClockHost.clockTr, Turing.ClockHost.hostCfg, Turing.ClockHost.hostCfg_symbols_yI, Turing.ClockHost.hostTM, Turing.ClockHost.runCfg, Turing.ClockHost.yard, Turing.ClockHost.yard_len, Turing.NDTM.stepWith, Turing.NDTM.stepWith_of_state, Turing.capOf}
\difficulty{3}
Let $S$ be a deterministic machine and $A$ a nondeterministic machine, both over the
Boolean alphabet, and fix a loop flag, an input word $x$, a natural number $m$, and a
choice bit $b$. For a configuration $a$ of $A$ on $x$, work-tape contents $sT$ and head
positions $sP$ for $S$, and a bit $v$, if $a$ is halted and its output is exactly the
one-symbol word $[v]$, then one step of the clock host from the run configuration built
from $a$, with yardstick head at position $m$, halts and has output $[v]$, preserving
the input position, both machines' work tapes and head positions, and the yardstick tape
with its head at $m$.
\end{theorem}

\begin{theorem}[Timeout failure transition]
\lean{Turing.ClockHost.timeout_fail}
\label{Turing.ClockHost.timeout_fail}
\leanok
\uses{Turing.Action.apply, Turing.Cfg, Turing.Cfg.inputSymbol, Turing.Cfg.workTapeSymbols, Turing.ClockHost.apply_hostCfg, Turing.ClockHost.clockTr, Turing.ClockHost.hostCfg, Turing.ClockHost.hostTM, Turing.ClockHost.idle, Turing.ClockHost.runCfg, Turing.ClockHost.slots, Turing.ClockHost.yI, Turing.ClockHost.yard, Turing.NDTM.stepWith, Turing.NDTM.stepWith_of_state, Turing.capOf, Turing.capOf_eq_some_some}
\difficulty{5}
Let $S$ be a finite-state deterministic Boolean machine and $A$ a nondeterministic
Boolean machine, and fix a Boolean loop parameter, Boolean input $x$, and natural number
$m$. Let $a$ be a configuration of $A$ on $x$ whose state is not absent with output a
singleton Boolean word. For any choice bit and any work-tape contents and head positions
for $S$, the host-machine step from the run configuration encoding $a$ at clock value
$m$ equals the following configuration: if looping is enabled, the host is in phase
$p_0$ with clock-tape head at $m-1$; otherwise the host is halted with clock-tape head
at $m$. In either case, the input position, work tapes, and their head positions are
inherited from $a$ and the given tapes for $S$, the clock tape contains the length-$m$
yardstick, and the output is empty.
\end{theorem}
